\subsection{Properties of Reduced Norms and Traces}

PROPOSITION 4. --- Let L be an extension of K and a an element of the central simple algebra A. We have the relations

\[
\operatorname{Pcrd} _ {\mathrm{A} _ {(\mathrm{L})} / \mathrm{L}} (1 \otimes a; \mathrm{X}) = \operatorname{Pcrd} _ {\mathrm{A} / \mathrm{K}} (a; \mathrm{X}),\tag{28}
\]

\[
\mathrm{Trd} _ {\mathrm{A} _ {(\mathrm{L})} / \mathrm{L}} (1 \otimes a) = \mathrm{Trd} _ {\mathrm{A} / \mathrm{K}} (a),\tag{29}
\]

\[
\mathrm{Nrd} _ {\mathrm{A} _ {(\mathrm{L})} / \mathrm{L}} (1 \otimes a) = \mathrm{Nrd} _ {\mathrm{A} / \mathrm{K}} (a).\tag{30}
\]

(``invariance under extension of scalars'').

The two sides of equality (28) have the same \(n\) -th power by the definition (formula (20) of VIII, p.~340) and the relation

\[
\mathrm{Pc} _ {\mathrm{A} _ {(\mathrm{L})} / \mathrm{L}} (1 \otimes a; \mathrm{X}) = \mathrm{Pc} _ {\mathrm{A} / \mathrm{K}} (a; \mathrm{X})
\]

(III, §9, No.~3, p.~544, formula (21)). Equality (28) therefore follows from Lemma 2 of VIII, p.~339. Formulas (29) and (30) follow from (28), (21), and (22).

COROLLARY 1. --- Let L be an extension of K. Let V be a vector space of dimension n over L and \(\theta\) a K-algebra homomorphism from A to \(\operatorname{End}_{\mathrm{L}}(\mathrm{V})\) . For every element a of A, we have

\[
\mathrm{Pcrd} _ {\mathrm{A/K}} (a; \mathrm{X}) = \chi_ {\theta (a)} (\mathrm{X}),\tag{31}
\]

\[
\mathrm{Trd} _ {\mathrm{A/K}} (a) = \mathrm{Tr} (\theta (a)),\tag{32}
\]

\[
\mathrm{Nrd} _ {\mathrm{A/K}} (a) = \det (\theta (a)).\tag{33}
\]

Let \(\theta'\) be the L-algebra homomorphism from \(\mathrm{A}_{(\mathrm{L})}\) to \(\mathrm{End}_{\mathrm{L}}(\mathrm{V})\) such that \(\theta'(\lambda \otimes a) = \lambda \theta(a)\) for \(\lambda \in \mathrm{L}\) and \(a \in \mathrm{A}\) . The algebra \(\mathrm{A}_{(\mathrm{L})}\) is simple by Corollary 2 of VIII, p.~222; the homomorphism \(\theta'\) is therefore injective. But the algebras \(\mathrm{A}_{(\mathrm{L})}\) and \(\mathrm{End}_{\mathrm{L}}(\mathrm{V})\) over the field L have the same degree, equal to \(n^2\) , so \(\theta'\) is an isomorphism. Corollary 1 then follows from Proposition 4 and Example 2 above.

COROLLARY 2. --- Let \(a\) and \(a'\) be elements of A and \(\lambda\) an element of K. We have the relations

\[
\mathrm{Pcrd} _ {\mathrm{A/K}} (a; a) = 0,\tag{34}
\]

\[
\mathrm{Pcrd} _ {\mathrm{A/K}} (\lambda a; \lambda \mathrm{X}) = \lambda^ {n} \mathrm{Pcrd} _ {\mathrm{A/K}} (a; \mathrm{X}),\tag{35}
\]

\[
\mathrm{Trd} _ {\mathrm{A/K}} (a + a ^ {\prime}) = \mathrm{Trd} _ {\mathrm{A/K}} (a) + \mathrm{Trd} _ {\mathrm{A/K}} (a ^ {\prime}),\tag{36}
\]

\[
\mathrm{Trd} _ {\mathrm{A/K}} (\lambda a) = \lambda \mathrm{Trd} _ {\mathrm{A/K}} (a),
\]

\[
\mathrm{Trd} _ {\mathrm{A/K}} (a a ^ {\prime}) = \mathrm{Trd} _ {\mathrm{A/K}} (a ^ {\prime} a),\tag{37}
\]

\[
\mathrm{Nrd} _ {\mathrm{A/K}} (a a ^ {\prime}) = \mathrm{Nrd} _ {\mathrm{A/K}} (a) \cdot \mathrm{Nrd} _ {\mathrm{A/K}} (a ^ {\prime}),\tag{38}
\]

\[
\mathrm{Nrd} _ {\mathrm{A/K}} (\lambda a) = \lambda^ {n} \mathrm{Nrd} _ {\mathrm{A/K}} (a),
\]

\[
\operatorname{Trd} _ {\mathrm{A} / \mathrm{K}} (1) = n, \operatorname{Nrd} _ {\mathrm{A} / \mathrm{K}} (1) = 1.\tag{39}
\]

Since A is central simple and has reduced degree n over K, there exist an extension L of K and a vector space V of dimension n over L such that \(\mathrm{A}_{(\mathrm{L})}\) is isomorphic to the algebra \(\mathrm{End}_{\mathrm{L}}(\mathrm{V})\) (VIII, p.~252, Theorem 1). Corollary 2 then follows from Corollary 1 and the properties of the trace and determinant of an endomorphism. In particular, formula (34) follows from the Cayley--Hamilton theorem (III, §8, No.~11, p.~541, Proposition 20, cf.~also VIII, p.~338, Remark 2).

COROLLARY 3. --- Let \(A^{o}\) be the opposite algebra of A. For every a in A, we have

\[
\mathrm{Pcrd} _ {\mathrm {A^ {\circ} /K}} (a; \mathrm{X}) = \mathrm{Pcrd} _ {\mathrm{A/K}} (a; \mathrm{X}),\tag{40}
\]

\[
\mathrm{Trd} _ {\mathrm {A^{\circ} /K}} (a) = \mathrm{Trd} _ {\mathrm{A/K}} (a),\tag{41}
\]

\[
\mathrm{Nrd} _ {\mathrm {A^{\circ} /K}} (a) = \mathrm{Nrd} _ {\mathrm{A/K}} (a).\tag{42}
\]

Choose an extension L of K, a vector space V of dimension n over L, and a homomorphism \(\theta\) from A to \(\operatorname{End}_{\mathrm{L}}(\mathrm{V})\) (VIII, p.~252, Theorem 1). Let \(V^{*}\) be the dual vector space of V. The mapping that sends an element a of A to the endomorphism \({}^{t}\theta(a)\) of \(V^{*}\) is a K-algebra homomorphism from \(A^{o}\) to \(\operatorname{End}_{\mathrm{L}}(\mathrm{V}^{*})\) . Corollary 3 then follows from Corollary 1 and Corollary 3 of III, §8, No.~4, p.~528.

The trace, norm, and characteristic polynomial of a are therefore the same whether we view a as an element of A or as an element of \(A^{\circ}\) . This property does not always hold when A is not assumed central simple (III, §9, p.~644, Exercise 1).

\pandocbounded{\includegraphics[keepaspectratio]{images/part002_d5c76b31.jpg}}

PROPOSITION 5. --- For any x in A, let \(t_{x}\) be the linear form \(y \mapsto \operatorname{Trd}_{\mathrm{A/K}}(xy)\) on A.

\begin{enumerate}
\def\labelenumi{\alph{enumi})}
\item
  The mapping \(t: x \mapsto t_x\) is an isomorphism of (A, A)-bimodules from A to its dual Hom\(_{K}\) (A, K).
\item
  Let \(h\) be a linear form on \(A\) . The following properties are equivalent:
\end{enumerate}

\begin{enumerate}
\def\labelenumi{(\roman{enumi})}
\item
  There exists an element \(\lambda\) of \(K\) such that \(h(x) = \lambda \operatorname{Trd}_{A / K}(x)\) for every \(x \in A\) .
\item
  We have \(h(xy) = h(yx)\) for all x, y in A.
\end{enumerate}

Recall (II, §1, No.~14, p.~225--226) that the (A, A)-bimodule structure on \(A^{*} = \operatorname{Hom}_{K}(A, K)\) is defined by the relation

\[
\langle a t b, c \rangle = \langle t, b c a \rangle\tag{43}
\]

for \(a, b, c \in A\) and \(t \in A^{*}\) . In particular, for every x in A, we have

\[
\begin{array}{c} \langle a t _ {x} b, c \rangle = \langle t _ {x}, b c a \rangle = \mathrm{Trd} _ {\mathrm{A/K}} (x b c a), \\ \langle t _ {a x b}, c \rangle = \mathrm{Trd} _ {\mathrm{A/K}} (a x b c), \end{array}
\]

and these two elements are equal by formula (37) of VIII, p.~342. We therefore have \(at_x b = t_{axb}\) , which means that \(t\) is a homomorphism of (A, A)-bimodules from A to A*.

We choose an extension L of the field K and an isomorphism \(\theta\) from the L-algebra \(\mathrm{A}_{(\mathrm{L})}\) to the matrix algebra \(\mathbf{M}_{n}(\mathrm{L})\) (VIII, p.~252, Theorem 1). We

identify the vector space \((\mathrm{A}^{*})_{(\mathrm{L})}\) with the dual of the vector space \(\mathrm{A}_{(\mathrm{L})}\) over the field L. By Proposition 4 of VIII, p.~341, with these conventions, we have

\[
\operatorname{Trd} _ {\mathrm{A} _ {(\mathrm{L})} / \mathrm{L}} = 1 _ {\mathrm{L}} \otimes \operatorname{Trd} _ {\mathrm{A} / \mathrm{K}}.\tag{44}
\]

Let \(t_{(L)}\) be the L-linear mapping from \(A_{(L)}\) to \(A_{(L)}^{*}\) deduced from t by extension of scalars; by formula (44) and Corollary 1, we have

\[
\langle t _ {(\mathrm{L})} (x), y \rangle = \operatorname{Trd} _ {\mathrm{A} _ {(\mathrm{L})} / \mathrm{L}} (x y) = \operatorname{Tr} (\theta (x) \theta (y))\tag{45}
\]

for \(x, y\) in \(\mathrm{A}_{(\mathrm{L})}\) . By Proposition 7 of II, §10, No.~11, p.~358, the mapping \(t_{(\mathrm{L})}\) is bijective; it follows that \(t\) is bijective. We have proved a).

Let \(h\) be in \(\mathrm{A}^*\) ; by the above, there exists an element \(a\) of \(\mathrm{A}\) such that \(h\) is equal to \(t_a\) . By a), we have

\[
h (x y) - h (y x) = t _ {a} (x y - y x) = t _ {a x} (y) - t _ {x a} (y).
\]

Consequently, the relation `` \(h(xy) = h(yx)\) for \(x,y\) in A'' is equivalent to `` \(t_{ax-xa} = 0\) for every \(x \in \mathrm{A}\) '', and by part a) of the proof, this means that \(a\) belongs to the center K of A. This proves b) in view of formula (36).

COROLLARY. --- The linear subspace of A generated by the elements of the form xy - yx, where x and y run through A, is a hyperplane, the kernel of the nonzero linear form \(Trd_{A/K}\) .

Remark. --- By formula (23) of VIII, p.~340, we have

\[
\mathrm{Tr} _ {\mathrm{A/K}} (a) = n \mathrm{Trd} _ {\mathrm{A/K}} (a)
\]

for every \(a \in A\) . If the characteristic of the field K is equal to 0 or a prime number p that does not divide n, then we can replace the reduced trace with the trace in Proposition 5. On the other hand, if the characteristic of K is a prime number that divides n, then we have \(\operatorname{Tr}_{\mathrm{A}/\mathrm{K}}(a) = 0\) for every \(a \in A\) .

